\subsection{Model Architecture}\label{subsec:architecture}

We first introduce the general structure of modern large language models (LLMs) to better demonstrate our approach and paradigm.

Modern LLMs involve Embedding, Transformer, and Output layers, as we denote them $\boldsymbol{E}$, $\boldsymbol{M}_n$, and $\boldsymbol{O}$ respectively, where $n$ indicates the layer of hidden state~\citep{DBLP:journals/corr/VaswaniSPUJGKP17, radford2019language}.

\subsection{Inference}\label{subsec:inference}

We used middle layers to enter the latent reasoning space to avoid breaking special usage in final layers~\citep{elhage2021mathematical}.

The whole inference procedure is described in (\ldots)

We first perform prefilling by the question with the length $q_0$ into the model, and get the prefilled KV Cache~\citep{pope2022efficientlyscalingtransformerinference} and the probability distribution of the answer's first token $t_{q_0+1}$ (starting to count from $t_1$).

After extracting the hidden state $\boldsymbol{h}_{n_0,q_0}$ (where $n_0$ denotes the layer number of the Transformer and $q_0$ denotes question length), we pass the state to a Variational Autoencoder~\citep{DBLP:journals/corr/abs-1906-02691}, and produce a vector for storage with length $l_s$.

Then, the new vector with the exact shape of $\boldsymbol{h}_{n,k}$, as we denote $\boldsymbol{h}^\prime_{n,k}$, will be operated by several gates, as described in Equation~\ref{eq:gate}.

\begin{equation}
    \boldsymbol{h}_{n,k} = \boldsymbol{h}_{n,k-1} \odot \boldsymbol{g} \odot \mathrm{DMG}(\boldsymbol{e}_k) \times
    s_\mathrm{inj}
    + \left(1-\boldsymbol{g}\right) \odot \boldsymbol{e}_k
    \label{eq:gate}
\end{equation}

where $\boldsymbol{e}_k=\mathrm{Embedding}(\boldsymbol{t}_k)$, where $\boldsymbol{t}_k$ denotes the token of sequence with position $k$, and $s_\mathrm{inj}$ denotes injection scalar, which is a constant scalar to bridge the difference in order of magnitude among gating value and the result of DMG (Dynamic Current-Step-Mixing Gate).

\paragraph{Gates} Our model has two gates: the general gate, the injection scalar, and the Dynamic Current Step Mixing Gate.

\begin{enumerate}
    \item \textbf{General Gate.} The general gate is derived from the sequence, its value is dynamic.
    We denote it as $\boldsymbol{g}_t$, and its derivation procedure can be described in Equation (not yet).
    We initialize it to all zeros, i.e., initially the model is equivalent to the original one, and it evolves as $\boldsymbol{g}$ moves away from $\boldsymbol{0}$.
    \item \textbf{Dynamic Current-Step-Mixing Gate (DMG).}
\end{enumerate}

\paragraph{Long-term Cache}

To enable long-term storage of reasoning features, we introduce a VAE-style bottleneck layer, compressing $1536$-dimensional features to $256$ dimensions and then expanding back to $1536$ before gating.

We store it in each operation, and concatenate these vectors into a matrix subsequently as a hidden state.

\subsection{Training}\label{subsec:training}

We use the DAPO objective as our reinforcement learning framework, but the whole procedure is slightly different from the original DAPO.

\subsubsection{Reinforcement Learning}

The whole procedure of reinforcement learning can be described in Figure.

We adapt the DAPO objective, which is formulated as:

\begin{equation}
    \begin{aligned}
        \mathcal{J}_{\mathrm{DAPO}}(\theta) =\quad& \mathbb{E}_{(q,a)\sim \mathcal{D}, \{o_i\}_{i=1}^G\sim \pi_{\theta_\mathrm{old}}(\cdot\mid q)}\\&
        \left[\frac{1}{\sum_{i=1}^{G}|o_i|}\sum_{i=1}^{G}\sum_{t=1}^{|o_i|}
        \min \left( r_{i,t}(\theta) \hat{A}_{i,t},
            \ \mathrm{clip} \left( r_{i,t}(\theta), 1 - {\varepsilon_{\mathrm{low}}}, 1 + {\varepsilon_{\mathrm{high}}} \right) \hat{A}_{i,t} \right) \right]
        \\
        \mathrm{subject\ to}\quad& 0< \left|\{o_i\mid\texttt{is\_equivalent}(a,o_i)\}\right|< G,
        \label{eq:dapoloss}
    \end{aligned}
\end{equation}
where
\begin{equation}
    r_{i,t}(\theta)=\frac{\pi_{\theta}(o_{i,t} \mid q, o_{i,<t})}{\pi_{\theta_{\mathrm{old}}}(o_{i,t} \mid q,o_{i,<t})},\quad\hat{A}_{i,t} = \frac{R_i - \mathrm{mean}(\{R_i\}_{i=1}^G)}{\mathrm{std}(\{R_i\}_{i=1}^G)}.
    \label{eq:advantage_calculation}
\end{equation}

Firstly we apply rigorous datasets to the inference procedure, and verify the answer~\citep{Kydlicek_Math-Verify_Math_Verification}.

We parallelized several inference procedures simultaneously, preferably $16$ for a same question, and verify the answer.

\paragraph{Award Mechanism}

We use $\mathrm{award}_{\mathrm{acc}}$ and $\mathrm{penality}_{\mathrm{acc}}$ to award or penalize in accuracy.
After mixing the accuracy and length evaluation, we inject it to loss.

\paragraph{Positive-Negative Ratio}

We use a ratio to mix positive and negative cases, which is $1:1$.

\subsection{Time Consistency Regularization}\label{subsec:time-consistency-reg}

A significant reason for chaotic response of multi-layer \textsc{Coconut} is its time inconsistency, which lead to the model's failure to maintain a consistent reasoning process.

We adapt a mechanism called ``Gating Value Bonus,'' which is used to guarantee the gate to converge to a fixed position, as we denote in Equation~\ref{eq:gating-value-bonus}.

\begin{equation}
    \mathcal{J}_{\mathrm{gate}}=\min\left\{C_0\exp\left[\lambda\left(\dfrac{1}{2}-g_i\right)^2\right]\right\}
    \label{eq:gating-value-bonus}
\end{equation}

\subsection{Parallelized Training}\label{subsec:parallelized-training}

After sampling, we concatenate the question, answer, and hidden state, and apply it to training procedure.